\subsection*{Encumbrance of armor}
Wearing armor is taxing and disruptive, especially on the limbs.
Armor pieces
- that are not INCONSPICUOUS or LIGHT -
cause penalties to its wearer depending on the location it covers:
\begin{itemize}
	\item[Head] -20 to Instinct and Charisma
	\item[Arms] -20 to Dex and Combat rolls
	\item[Legs] -20 to Agility, -2 MS
\end{itemize}

\subsection*{A note on "reusing" armor}
When a character acquires one or more pieces of armor from fallen combatants,
multiple stipulations apply:

\paragraph{Armor generally does not fit.}
If an armor piece is not bought,
there is a big chance for it not to fit properly.
\\%
Roll a D5.
On a 3, the armor fits well.
\\%
On a higher result, the armor is too big.
While it can be worn,
all penalties are doubled from the encumbrance.
\\%
On a lower result, the armor is too small.
The character simply won't fit
- it's scrap at best.
\\%
If the characters are of different Sizes (see p. \pageref{sizestable}),
there is no chance for the armor to fit.

\paragraph{Armor gets damaged.}
Most armor protects for only a few hits.
If the armor blocked multiple hits,
it is likely useless now.
\\%
Roll a D5.
If the armor was hit that many times,
it can only be used as scrap.
\\%
If the armor blocked incredibly heavy hits,
the GM may obviously forego the roll.

\paragraph{PES have to be calibrated.}
Personal energy shields,
as encase the wearer very closely,
have to be calibrated to its user.
This takes technical knowledge (or the instruction manual)
and about 20 minutes.
\\%
They also require a L1NC to control,
meaning anyone without the implant cannot calibrate and activate a PES,
and a PES may be bio-coded.
